% ============================================================
% Ancienne version de la section 2.2 « Fonctions booléennes à seuil (SBF) »
% (juillet 2026, passages de Claude repris en partie par Quentin), remplacée le
% 2026-10-04 par une réécriture complète demandée par Quentin. Non compilé.
%
% Contenu retiré du tableau des échelles, à reprendre ailleurs si besoin :
% - colonne PTBN = |T_n|^n - (|T_n|-1)^n (dynamiques ayant au moins un nœud de contrôle,
%   donc un vecteur non trivial en PTBN) : section « Quel paysage explore-t-on ? », ou MCSBN
%   (1 - PTBN/SBN = 98,9 % en d = 4, à rapprocher des 98,89 % de vecteurs triviaux de MCSBN) ;
% - colonnes vec. (|H_n| = 2, 4, 10, 36, 202, 1828, partitions binaires de 2^n), btree
%   (A_n = 1 + A_{n-1}^2) et tagged (T_n = 1 + n T_{n-1}^2) : section Correction, où les
%   programmes setup_btree, treeinsolver et taggedtree sont présentés ;
% - colonne Poids = (2B+1)^{n^2} : section Statistiques, si l'on parle de l'espace des génotypes.
% Valeurs revérifiées le 2026-10-04 (force brute sur les poids et w_bound_orbits.py) :
% toutes exactes. Défauts relevés : « 2^n hyperplans » (il y en a 2^n - 1, x ≠ 0) ;
% comparaison des nombres d'arbres au nombre de SBN (« vec., btree, tagged << SBN »), qui
% ne mesurent pas la même chose ; notation du code w_bound(n).
% ============================================================

\subsection{Fonctions booléennes à seuil (SBF)}
\label{sec:sbf}

[Cf. rapport Ailing]

Commentaire sur les permutations des poids des SBF et la non prise en compte d'une arête non existante (chemin à exploiter). Ceci explique les différences entre cette table d'échelle et celles qui ont étés faits précédemment. 

\begin{table}[H]
\centering
\small
\setlength{\tabcolsep}{4pt}
\caption{Ordres de grandeur par dimension~$n$. \textbf{SBF}~$=|\mathcal{T}_n|$ (chambres,
BFS exact)~; \textbf{orb.}~$=|\mathcal{T}_n/S_n|$ (réduction par permutation des nœuds)~;
\textbf{SBN}~$=|\mathcal{T}_n|^n$ (\emph{produit de SBF}~: tout $n$-uplet de SBF est
réalisable, cf. indépendance des colonnes)~; \textbf{PTBN}~$=|\mathcal{T}_n|^n-(|\mathcal{T}_n|-1)^n$
(SBN admettant $\ge 1$ nœud de contrôle, i.e. une colonne $f_c=\pi_c$~; forme close validée par
énumération)~; \textbf{$B$}~$=w_{\text{bound}}(n)$ (ILP, pire cas)~; \textbf{Poids}~$=(2B+1)^{n^2}$
(espace syntaxique des matrices)~; \textbf{vec.} vecteurs de décomposition (couverts par
\texttt{setup\_btree})~; \textbf{btree} arbres non étiquetés $A(d)=1+A(d-1)^2$
(\texttt{treeinsolver})~; \textbf{tagged} arbres étiquetés par nœud de contrôle
$T(m)=1+m\,T(m-1)^2$ (\texttt{taggedtree}).}
\label{tab:echelles}
\begin{tabular}{r r r r r c r r r r}
\toprule
$n$ & SBF & orb. & SBN & PTBN & $B$ & Poids & vec. & btree & tagged \\
\midrule
\emph{form.} & $|\mathcal{T}_n|$ & $|\mathcal{T}_n/S_n|$ & $|\mathcal{T}_n|^n$ & $|\mathcal{T}_n|^n{-}(|\mathcal{T}_n|{-}1)^n$ & $w_{\mathrm{b}}(n)$ & $(2B{+}1)^{n^2}$ & $|H_n|$ & $1{+}A_{n-1}^2$ & $1{+}n\,T_{n-1}^2$ \\
\midrule
1 & 2      & 2   & $2$                 & $1$                 & 1 & $3$                 & 2   & 2                 & 2 \\
2 & 6      & 4   & $36$                & $11$                & 1 & $81$                & 4   & 5                 & 9 \\
3 & 32     & 10  & $3.28{\cdot}10^{4}$ & $2977$              & 2 & $1.95{\cdot}10^{6}$ & 10  & 26                & 244 \\
4 & 370    & 36  & $1.87{\cdot}10^{10}$& $2.02{\cdot}10^{8}$ & 3 & $3.32{\cdot}10^{13}$& 36  & 677               & $2.38{\cdot}10^{5}$ \\
5 & 11292  & 224 & $1.84{\cdot}10^{20}$& $8.13{\cdot}10^{16}$& 5 & $1.08{\cdot}10^{26}$& 202 & $4.58{\cdot}10^{5}$ & $2.84{\cdot}10^{11}$ \\
6 & $1.07{\cdot}10^{6}$ & 3024 & $1.47{\cdot}10^{36}$& $8.26{\cdot}10^{30}$& 9 & $1.08{\cdot}10^{46}$& 1828 & $2.10{\cdot}10^{11}$ & $4.82{\cdot}10^{23}$ \\
\bottomrule
\end{tabular}
\end{table}

Les colonnes \textbf{Poids} et \textbf{SBN} explosent (respectivement $(2B+1)^{n^2}$ et
$|\mathcal{T}_n|^n$)~: la force brute est intenable dès $n=4$. Les trois dernières colonnes
donnent les trois leviers qui rendent la montée en dimension possible~: réduction par $S_n$
(facteur jusqu'à $n!$), restriction au sous-espace pluripotent (\emph{vec.}, \emph{btree},
\emph{tagged} $\ll$ SBN), et énumération par arbre-dans-le-solveur plutôt que par grille de
poids. Les trois familles d'arbres sont emboîtées, $\text{vec.}\le\text{btree}\le\text{tagged}$,
et correspondent à trois outils de complétude croissante~: \texttt{setup\_btree} (une orientation
par vecteur, \emph{incomplet}) $\subset$ \texttt{treeinsolver} (toutes orientations) $\subset$
\texttt{taggedtree} (étiqueté, sans permutation).

\paragraph{Borne des poids par programmation linéaire en nombre entier.}

On note $\mathcal{T}_n = \{ f_w : w \in \mathbb{Z}^n \}$ l'ensemble des SBF réalisables
en dimension~$n$.

Pour une SBF réalisable $f$, posons $P = \{\mathbf{x} : f(\mathbf{x}) = 1\}$ et
$N = \{\mathbf{x} : f(\mathbf{x}) = 0\}$. La borne de poids minimale réalisant $f$ est
$W(f) = \min\{ \|w\|_\infty : w \in \mathbb{Z}^n,\ f_w = f \}$. Comme
$\sum_i w_i x_i \in \mathbb{Z}$ pour $w$ entier, la condition stricte
$\sum_i w_i x_i > 0$ équivaut à $\sum_i w_i x_i \ge 1$, ce qui permet d'écrire le problème comme un programme linéaire en nombres entiers:
\begin{equation}
\begin{aligned}
W(f) = \min_{w \in \mathbb{Z}^n,\ B \in \mathbb{Z}_{\ge 0}}\quad & B\\
\text{s.c.}\quad
& \textstyle\sum_i w_i x_i \ge 1 && \forall\, \mathbf{x} \in P,\\
& \textstyle\sum_i w_i x_i \le 0 && \forall\, \mathbf{x} \in N,\\
& -B \le w_i \le B && \forall\, i \in \{1,\dots,n\}.
\end{aligned}
\end{equation}
Les deux premières familles de contraintes caractérisent $f$ ~;  $-B \le w_i \le B$
linéarise $B = \|w\|_\infty$.

La borne recherchée, celle qui permet d'exprimer \emph{toutes} les SBF de
dimension~$n$, est alors le pire cas sur $\mathcal{T}_n$~:
\begin{equation}
\boxed{\max_{f_w \in \mathcal{T}_n} W(f)\ }
\end{equation}
c'est-à-dire le plus petit $B$ tel que tout $f_w \in \mathcal{T}_n$ soit réalisable avec
$\|w\|_\infty \le B$.

Il ne reste donc plus qu'à énumérer toutes les SBF et résoudre l'ILP ci-dessus pour chacune. Enumérer les $2^{2^n}$ fonctions booléennes pour ne selectionner que les réalisables n'est pas raisonnable pour $n\ge 5$, il faut donc trouver une autre manière. Une idée est de passer par les poids, il suffirait de trouver \emph{une} borne minimale des poids garantissant l'obtention de toutes les SBF, mais il n'en a pas été trouvé de suffisamment raisonnable pour le calcul. L'idée retenue a donc été d'explorer l'espace des SBF par un parcourt (en largeur) de graphe, en considérant que deux SBF sont voisines si on peut passer de l'une à l'autre par changement d'un unique bit, i.e. en traversant un unique hyper-plan de $\mathbb{R}^n$ dans l'espace des poids. [Cf la thèse de Rémi, les parties 3 et 4 présentent une approche différente d'énumération que celle évoquée ici, mais ses illustrations permettent d'y voir plus clair]

\paragraph{Complétude du parcours.} Cette relation de voisinage est exactement l'adjacence des \emph{chambres} de l'arrangement des $2^n$ hyperplans $H_{\mathbf x} = \{w : \sum_i w_i x_i = 0\}$ dans l'espace des poids~: sur chaque chambre (région ouverte délimitée par ces hyperplans), le signe de $\sum_i w_i x_i$ est constant pour tout $\mathbf x$, donc $f_w$ y est constante. $\mathcal{T}_n$ est donc exactement l'ensemble des chambres, et deux chambres voisines (séparées par une seule facette) donnent deux SBF qui ne diffèrent que d'un bit.

Le graphe des chambres est connexe~: $\mathbb{R}^n$ (l'espace complet, avant de retirer les hyperplans) est connexe, donc deux poids quelconques $w,w'$ sont reliés par un chemin~; quitte à le perturber légèrement, ce chemin ne traverse jamais deux hyperplans au même instant (cas de mesure nulle), donc il les traverse un par un, chaque traversée faisant passer à une chambre voisine. La suite des chambres ainsi traversées relie celle de $w$ à celle de $w'$~: le graphe est connexe.

Un parcours en largeur partant d'une seule SBF réalisable (par exemple $f=0$, obtenue avec $w=0$) et ne testant, à chaque étape, que la réalisabilité du voisin à un bit de distance (simple faisabilité LP) visite donc nécessairement toute $\mathcal{T}_n$~: la connexité garantit qu'aucune SBF n'est hors d'atteinte, contrairement à une recherche par force brute sur les poids où une SBF exigeant de gros poids peut être manquée sans avertissement.

\paragraph{Réduction par symétrie.} Deux SBF liées par une permutation des $n$ coordonnées ont la même borne $\|w\|_\infty$ (permuter les poids ne change pas leur norme), il suffit donc d'explorer un seul représentant par orbite sous $S_n$ . Cela ne coûte rien en complétude~: si $r$ est le représentant exploré et $\sigma \cdot r$ un autre représentant de la même orbite, retourner le bit $k$ de $\sigma\cdot r$ revient, après application de $\sigma^{-1}$, à retourner le bit $\sigma^{-1}(k)$ de $r$~; comme $\sigma^{-1}$ est une bijection sur les positions, les $2^n$ voisins de $r$ (une fois ramenés à leur représentant canonique) donnent déjà toutes les orbites voisines, quel que soit le représentant de départ choisi.


\begin{table}[H]
\centering
\caption{Les 6 SBF distinctes en dimension $2$. $f(\mathbf{x}) = 1$ ssi $\sum_i w_i x_i > 0$, avec $w_i \in \{-1, 0, 1\}$.}
\label{tab:sbf2}
\begin{tabular}{cccccl}
\toprule
$(w_0,\,w_1)$ repr. & $f(0,0)$ & $f(1,0)$ & $f(0,1)$ & $f(1,1)$ & Expression \\
\midrule
$(-1,-1)$ & 0 & 0 & 0 & 0 & $0$ \\
$(1,-1)$  & 0 & 1 & 0 & 0 & $x_0 \wedge \neg x_1$ \\
$(-1,\;1)$  & 0 & 0 & 1 & 0 & $\neg x_0 \wedge x_1$ \\
$(1,\;0)$   & 0 & 1 & 0 & 1 & $x_0$ \\
$(0,\;1)$   & 0 & 0 & 1 & 1 & $x_1$ \\
$(1,\;1)$   & 0 & 1 & 1 & 1 & $x_0 \vee x_1$ \\
\bottomrule
\end{tabular}
\end{table}

